% from ega1-8-fr.tex
\oldpage[I]{168}
une immersion locale (resp. un isomorphisme local), car la question étant
locale sur $X$ et $Y$, on peut se borner pour le démontrer au cas où $f$ est
une immersion fermée (resp. une immersion ouverte) et notre assertion découle
alors aussitôt de (\hyperref[I.4.2.2-fr]{4.2.2}).

On peut donc supposer que $f$ est un morphisme \emph{dominant}, ce qui
entraîne que $Y$ est, lui aussi, irréductible
(\hyperref[0.2.1.5-fr]{0, 2.1.5}), donc que $X$ et $Y$ sont tous deux
\emph{intègres}. En outre, la question étant locale sur $Y$, on peut supposer
que $Y$ est un schéma affine~; comme $f$ est séparé, $X$ est un schéma
(\hyperref[I.5.5.1-fr]{5.5.1, (ii)}), et on est finalement dans les
hypothèses de (\hyperref[I.8.2.7-fr]{8.2.7}). Alors, pour tout $x\in X$,
$\theta_x^\sharp$ est injectif~; mais l'hypothèse que $f$ est une immersion
locale implique que $\theta_x^\sharp$ est surjectif
(\hyperref[I.4.2.2-fr]{4.2.2}), donc $\theta_x^\sharp$ est bijectif,
autrement dit (avec l'identification de
(\hyperref[I.8.2.7-fr]{8.2.7})) on a
$\mathscr{O}_{\psi(x)}=\mathscr{O}_x$. Cela implique par
(\hyperref[I.8.2.4-fr]{8.2.4}) que $\psi$ est une application
\emph{injective}, ce qui démontre déjà la proposition lorsque $f$ est un
isomorphisme local (\hyperref[I.4.5.3-fr]{4.5.3}). Lorsqu'on suppose seulement
que $f$ est une immersion locale, pour tout $x\in X$ il existe un voisinage
ouvert $U$ de $x$ dans $X$ et un voisinage ouvert $V$ de $\psi(x)$ dans $Y$
tels que la restriction de $\psi$ à $U$ soit un homéomorphisme de $U$ sur une
partie \emph{fermée} de $V$. Or, $U$ est dense dans $X$, donc $\psi(U)$ est
dense dans $Y$ et \emph{a fortiori} dans $V$, ce qui prouve que
$\psi(U)=V$~; comme $\psi$ est injective, $\psi^{-1}(V)=U$ et ceci achève de
montrer que $\psi$ est un homéomorphisme de $X$ sur $\psi(X)$.

\subsection{Les schémas de Chevalley.}
\label{subsection:I.8.3-fr}

\begin{env}[8.3.1]
\label{I.8.3.1-fr}
Soient $X$ un schéma intègre \emph{noethérien}, $R$ son corps des fonctions
rationnelles~; désignons par $X'$ l'ensemble des sous-anneaux locaux
$\mathscr{O}_x\subset R$, où $x$ parcourt $X$. L'ensemble $X'$ vérifie les
trois conditions suivantes~:
\begin{enumerate}
  \item[(Sch. 1)] \emph{Pour tout $M\in X'$, $R$ est le corps des fractions
    de $M$.}
  \item[(Sch. 2)] \emph{Il existe un ensemble fini de sous-anneaux
    noethériens $A_i$ de $R$ tels que $X'=\bigcup_i L(A_i)$ et que, pour tout
    couple d'indices $i$, $j$, le sous-anneau $A_{ij}$ de $R$ engendré par
    $A_i\cup A_j$ soit une algèbre de type fini sur $A_i$.}
  \item[(Sch. 3)] \emph{Deux éléments $M$, $N$ de $X'$ qui sont apparentés
    sont identiques.}
\end{enumerate}
\end{env}

On a en effet vu dans (\hyperref[I.8.2.1-fr]{8.2.1}) que (Sch. 1) est
satisfaite, et (Sch. 3) résulte de (\hyperref[I.8.2.2-fr]{8.2.2}). Pour
démontrer (Sch. 2), il suffit de recouvrir $X$ par un nombre fini d'ouverts
affines $U_i$, d'anneaux noethériens, et de prendre
$A_i=\Gamma(U_i,\mathscr{O}_X)$~; l'hypothèse que $X$ est un schéma entraîne
que $U_i\cap U_j$ est affine et que
$\Gamma(U_i\cap U_j,\mathscr{O}_X)=A_{ij}$
(\hyperref[I.5.5.6-fr]{5.5.6})~; en outre, comme l'espace $U_i$ est
noethérien, l'immersion $U_i\cap U_j\to U_i$ est de type fini
(\hyperref[I.6.3.5-fr]{6.3.5}), donc $A_{ij}$ est une $A_i$-algèbre de type
fini (\hyperref[I.6.3.3-fr]{6.3.3}).

\begin{env}[8.3.2]
\label{I.8.3.2-fr}
Les structures dont les axiomes sont (Sch. 1), (Sch. 2) et (Sch. 3)
généralisent les «~schémas~» au sens de C.~Chevalley, qui suppose en outre que
$R$ est une extension de type fini d'un corps $K$ et que les $A_i$ sont des
$K$-algèbres de type fini (ce qui rend inutile une partie de (Sch. 2))
\cite{I-1}. Inversement, si on a une telle structure sur un ensemble $X'$,
on peut lui associer un schéma intègre $X$ en utilisant les remarques de
(\hyperref[I.8.2.6-fr]{8.2.6})~: l'espace sous-jacent de $X$ est égal à $X'$
muni de la topologie définie dans (\hyperref[I.8.2.6-fr]{8.2.6}), et du
faisceau $\mathscr{O}_X$

